
% ---------------------------------------------------------------------------------------------
% NOTE. This file currently holds only the two sections moved out of the Discussion chapter:
% Limitations and Future Work. The rest of the Conclusion (introduction, summary of findings,
% contributions, closing remarks) still needs to be written around them.
%
% Both sections were trimmed when they moved. The parts that INTERPRET a finding now live in
% Section~\ref{sec:disc:limits}; what remains here is the caveat each places on a claim. Check
% against that section before adding material back, or the two will restate each other.
% ---------------------------------------------------------------------------------------------

\chapter{Conclusion}\label{ch:conclusion}

\section{Limitations}\label{sec:conc:limitations}

\paragraph{Simulation only.}
Every result is from SafeLIBERO. No physical robot was used and no claim of sim-to-real transfer is
made.

\paragraph{No formal guarantee is inherited.}
The distilled policy carries no runtime constraint, so nothing of the shield's formal argument
survives its removal. The conditions under which a learned controller instead inherits
input-to-state safety~\cite{cosner2022il_cbf} are set out in Section~\ref{sec:bg:learned}, and none
holds here: the shield is a nominal rather than a robust QP, the plant is not a known
control-affine system, and neither a bounded learning error nor a Lipschitz constant is available
for a 3-billion-parameter flow-matching policy. What is demonstrated is an empirical reduction in
collisions.

\paragraph{Privileged geometry.}
The barrier is constructed from obstacle poses and meshes read directly from the simulator. A
deployed system would estimate that geometry from perception and inherit the resulting error, so
every shielded figure reported here is an upper bound for this class of filter. The distilled
policy uses no geometry at inference, but the dependency remains during \emph{training}.

\paragraph{Approximate robot model.}
The barrier represents the robot with spheres rather than its exact shape, because the exact
version would not solve fast enough to run at every control step (Section~\ref{sec:barrier}). How
accurate that approximation actually was is never measured directly; the account in
Section~\ref{sec:disc:limits} of where the transfer stops is a reading of which bodies collided,
not a measurement of constraint fidelity.

\paragraph{Single training seed.}
Each distilled policy is one training run and seed variance is unquantified.

\paragraph{Single policy and action space.}
Everything here uses one model, $\pi_{0.5}$, which produces chunks of actions by integrating a
flow field. Two results depend on that design --- the reinforcement-learning negative and the
best-of-$N$ measurement --- and may not carry over to policies that emit one action at a time.

\paragraph{Scope of the negative results.}
Each negative is a finding about the implementation and budget tested, not about its method class.
The reinforcement-learning result concerns the reward design and the variant that could be run at
this scale; richer constraint signals exist and were out of computational reach
(Section~\ref{sec:disc:limits}). The language result concerns a single phrasing. The best-of-$N$
measurement supports a state-dominant account of collision risk in this policy and benchmark rather
than a general claim about generative policies.

\paragraph{Attribution caveats.}
The per-body attribution is a documented lower bound that can miss delayed or indirect contacts.

\section{Future Work}\label{sec:conc:future}

\textbf{A physical robot} is the most direct extension. Whether internalised avoidance survives
contact dynamics, calibration error and a real perception pipeline is the question this work most
needs answered, and Section~\ref{sec:disc:cost} suggests the deployment configuration to test: the
distilled policy with the shield reattached, which reached the highest collision-avoidance rate of
any arm evaluated while keeping the runtime constraint a safety case would require.

\textbf{Repeating the training across seeds} is the cheapest outstanding item and addresses the
weakest one. Each distilled policy here is a single run, so the training variance behind the
headline is unquantified even though the evaluation variance is not.

\textbf{Perception-grounded barriers} would remove the privileged-geometry dependency and measure
how much of the reported safety survives estimation error --- the same question a physical
deployment asks, isolated from the rest of the reality gap. Learned safety representations offer
one route~\cite{meng2024conbat,sun2025lpb}, though both retain the inference-time correction this
work removes.

\textbf{A whole-body barrier} with proper link geometry and directly constrained link velocities
would lower the shield's own floor (Section~\ref{sec:res:scope}); whether that improvement then
distils is separately testable. It would also make CBF-guided sampling worth testing, since the
existing implementation was verified against a closed-form projection but wired to an
end-effector barrier, and so never activated on the scenes where arm links collided.

\textbf{Other architectures and moving obstacles} bound the generality of the result. The
best-of-$N$ and reinforcement-learning findings depend on chunked flow sampling and may not hold for
policies emitting single actions, and every scene here is static --- the case that matters for
deployment is a constraint that moves.

Finally, \textbf{separating the surviving explanations} would say something general about when
demonstrations transfer between controllers. Observation aliasing was measured and rejected
(Section~\ref{sec:res:source}) and state coverage is supported but measured only on
end-effector position; distribution shift and action-distribution mismatch remain. A direct
comparison against model-based safe RL on this benchmark~\cite{safedojo} would place distillation
against the strongest current alternative.
